\section{Introduction}
\label{sec:introduction}

State the problem and explain why it matters. Give the reader the minimum
background needed to understand the rest of the paper, then summarize the
central idea and contributions. Keep this section focused on the paper's
claims rather than on implementation details.

This repository follows the official USENIX conference paper template
\cite{usenix-paper-templates}. The template uses US letter paper, two columns,
a 7 by 9 inch text block, and 10-point Times-style text with 12-point leading.
The conference call for papers remains the authority for page limits,
anonymity, and supplemental-material rules.

A useful contribution list is:
\begin{itemize}
  \item The first concrete contribution, stated as a result or technique.
  \item The second contribution, with its scope and assumptions.
  \item The evaluation or artifact contribution and the evidence supporting it.
\end{itemize}

The rest of the paper is organized as follows. Section~\ref{sec:method}
describes the approach, Section~\ref{sec:evaluation} evaluates it, and
Section~\ref{sec:related-work} places it in context.
